% =====================================================================
\chapter{Técnicas de integración}
\label{tema:3}
% =====================================================================

La regla de Barrow reduce el cálculo de $\int_a^b f$ a encontrar una primitiva de $f$. A diferencia de la derivación, que sigue reglas mecánicas, la búsqueda de primitivas es un arte con unas pocas herramientas: el cambio de variable, que invierte la regla de la cadena; la integración por partes, que invierte la regla del producto; la descomposición en fracciones simples, que resuelve por completo las funciones racionales; y las identidades trigonométricas. Algunas funciones elementales, como $e^{-x^{2}}$, no tienen primitiva elemental, y para ellas se recurre a las funciones definidas por integrales del \cref{tema:2} o a las series del \cref{tema:6}. Se siguen el capítulo~19 de Spivak~\cite{spivak}, las §5.5 y~5.6 de~\cite{openstax1} y §3.1 a~3.5 de~\cite{openstax2}, el capítulo~7 de Strang~\cite{strang} y la unidad~4 del curso 18.01SC~\cite{mit1801}.

\section{Cambio de variable}

\begin{teorema}[Cambio de variable]\label{teo:sustitucion}
Sea $\varphi\colon[\alpha,\beta]\to\R$ con derivada continua, y sea $f$ continua en un intervalo $I$ que contiene a $\varphi([\alpha,\beta])$. Entonces
\[
\int_{\alpha}^{\beta}f\bigl(\varphi(t)\bigr)\,\varphi'(t)\dif t=\int_{\varphi(\alpha)}^{\varphi(\beta)}f(x)\dif x .
\]
En forma indefinida: si $F'=f$, entonces $\int f(\varphi(t))\,\varphi'(t)\dif t=F(\varphi(t))+C$.
\end{teorema}

\begin{proof}
Por el \cref{cor:existencia-primitiva}, $f$ tiene una primitiva $F$ en $I$. Por la regla de la cadena, $(F\circ\varphi)'(t)=f(\varphi(t))\,\varphi'(t)$, que es continua, y la regla de Barrow aplicada a ambos miembros da
$\int_\alpha^\beta f(\varphi(t))\varphi'(t)\dif t=F(\varphi(\beta))-F(\varphi(\alpha))=\int_{\varphi(\alpha)}^{\varphi(\beta)}f(x)\dif x$.
\end{proof}

En la práctica se escribe $u=\varphi(t)$, $\dif u=\varphi'(t)\dif t$, y se sustituye formalmente. En una integral definida conviene cambiar también los límites, lo que evita deshacer el cambio.

\begin{ejemplo}\label{ej:sustitucion}
(a) $\int x\,e^{x^{2}}\dif x$: con $u=x^{2}$, $\dif u=2x\dif x$, queda $\frac12\int e^{u}\dif u=\frac12e^{x^{2}}+C$.
(b) $\int\tan x\dif x=\int\frac{\sin x}{\cos x}\dif x$: con $u=\cos x$, $\dif u=-\sin x\dif x$, queda $-\int\frac{\dif u}{u}=-\ln\abs{\cos x}+C$. En general, $\int\frac{g'(x)}{g(x)}\dif x=\ln\abs{g(x)}+C$.
(c) $\int_0^1\frac{x}{1+x^{2}}\dif x$: con $u=1+x^{2}$, los límites pasan a $u=1$ y $u=2$, y la integral vale $\frac12\int_1^2\frac{\dif u}{u}=\frac{\ln2}{2}$.
\end{ejemplo}

El teorema también se usa en sentido inverso: para calcular $\int_a^b f(x)\dif x$ se escribe $x=\varphi(t)$ con $\varphi$ biyectiva de $[\alpha,\beta]$ en $[a,b]$, $\varphi(\alpha)=a$, $\varphi(\beta)=b$, y se integra en $t$.

\begin{ejemplo}[Área del círculo]\label{ej:circulo}
Con $x=\sin t$, $t\in[-\pi/2,\pi/2]$, se tiene $\sqrt{1-x^{2}}=\cos t$ (pues $\cos t\ge0$ ahí) y $\dif x=\cos t\dif t$:
\[
\int_{-1}^{1}\sqrt{1-x^{2}}\dif x=\int_{-\pi/2}^{\pi/2}\cos^{2}t\dif t=\int_{-\pi/2}^{\pi/2}\frac{1+\cos2t}{2}\dif t=\frac\pi2 .
\]
Es el área del semicírculo de radio~1; el círculo completo tiene área $\pi$.
\end{ejemplo}

\section{Integración por partes}

\begin{teorema}[Integración por partes]\label{teo:partes}
Sean $u,v$ con derivada continua en $[a,b]$. Entonces
\[
\int_a^b u(x)\,v'(x)\dif x=\eval{u(x)\,v(x)}{a}{b}-\int_a^b u'(x)\,v(x)\dif x ,
\]
y en forma indefinida, $\int u\,v'=u\,v-\int u'\,v$, o abreviadamente $\int u\dif v=uv-\int v\dif u$.
\end{teorema}

\begin{proof}
Por la regla del producto, $(uv)'=u'v+uv'$, función continua. La regla de Barrow da $\int_a^b(u'v+uv')=\eval{uv}{a}{b}$, y basta despejar usando la linealidad.
\end{proof}

La técnica es útil cuando $u$ se simplifica al derivarla (un polinomio, un logaritmo, una función trigonométrica inversa) y $v'$ no se complica al integrarla (una exponencial, un seno, un coseno).

\begin{ejemplo}\label{ej:partes}
(a) $\int x\,e^{x}\dif x$: con $u=x$, $v'=e^{x}$, se obtiene $x\,e^{x}-\int e^{x}\dif x=(x-1)e^{x}+C$.
(b) $\int\ln x\dif x$: con $u=\ln x$, $v'=1$, se obtiene $x\ln x-\int x\cdot\frac1x\dif x=x\ln x-x+C$.
(c) $J=\int e^{x}\sin x\dif x$. Dos integraciones por partes, con $u=\sin x$ y después $u=\cos x$, dan $J=e^{x}\sin x-\int e^{x}\cos x\dif x=e^{x}\sin x-e^{x}\cos x-J$. Despejando, $J=\frac12e^{x}(\sin x-\cos x)+C$. La integral reaparece y se despeja: es una integral \emph{cíclica}.
\end{ejemplo}

La integración por partes repetida conduce a \emph{fórmulas de reducción}, que expresan una integral con un parámetro $n$ en función de la misma integral con $n-1$ o $n-2$ y se resuelven por inducción (\cref{eje:3-reduccion,eje:3-wallis}).

\section{Funciones racionales y fracciones simples}

Una función racional $P/Q$ es \emph{propia} si $\operatorname{gr}P<\operatorname{gr}Q$. Toda función racional es un polinomio más una función racional propia, por la división de polinomios. Todo polinomio real se factoriza en factores lineales $(x-r)$ y cuadráticos irreducibles $(x^{2}+px+q)$, con $p^{2}-4q<0$.

\begin{teorema}[Descomposición en fracciones simples]\label{teo:fracciones}
Sea $P/Q$ una función racional propia y supóngase $Q$ mónico factorizado como producto de potencias de factores lineales $(x-r)^{k}$ y cuadráticos irreducibles $(x^{2}+px+q)^{m}$. Entonces $P/Q$ es, de forma única, una suma de términos
\[
\frac{A_j}{(x-r)^{j}}\ \ (1\le j\le k)\qquad\text{y}\qquad\frac{B_j\,x+C_j}{(x^{2}+px+q)^{j}}\ \ (1\le j\le m),
\]
uno por cada potencia de cada factor.
\end{teorema}

La demostración general es algebraica~\cite[cap.~19]{spivak}. El caso más frecuente, con raíces reales simples, admite una prueba breve que además da los coeficientes.

\begin{proposicion}[Raíces simples]\label{prop:raices-simples}
Sea $Q(x)=(x-r_1)\cdots(x-r_n)$ con $r_1,\dots,r_n$ distintos y sea $P$ un polinomio de grado menor que $n$. Entonces
\[
\frac{P(x)}{Q(x)}=\sum_{i=1}^{n}\frac{A_i}{x-r_i},\qquad A_i=\frac{P(r_i)}{\prod_{j\neq i}(r_i-r_j)}=\frac{P(r_i)}{Q'(r_i)} .
\]
\end{proposicion}

\begin{proof}
Sea $R(x)=\sum_{i=1}^{n}A_i\prod_{j\neq i}(x-r_j)$, con los $A_i$ del enunciado. $R$ es un polinomio de grado menor que $n$. En $x=r_k$ todos los productos con $i\neq k$ contienen el factor $(r_k-r_k)=0$, así que $R(r_k)=A_k\prod_{j\neq k}(r_k-r_j)=P(r_k)$. Los polinomios $R$ y $P$, de grado menor que $n$, coinciden en $n$ puntos distintos, luego $R-P$ tiene $n$ raíces y grado menor que $n$, y es el polinomio nulo. Dividiendo $R=P$ por $Q$ se obtiene la igualdad. Por último, derivando $Q$ como producto, todos los sumandos de $Q'(r_i)$ se anulan salvo $\prod_{j\neq i}(r_i-r_j)$.
\end{proof}

Las fracciones simples tienen primitivas conocidas: $\int\frac{A\dif x}{x-r}=A\ln\abs{x-r}+C$; $\int\frac{A\dif x}{(x-r)^{j}}=-\frac{A}{(j-1)(x-r)^{j-1}}+C$ si $j\ge2$; y para el factor cuadrático se completa el cuadrado, $x^{2}+px+q=(x+\frac p2)^{2}+s^{2}$ con $s^{2}=q-\frac{p^{2}}{4}>0$, y se separa la parte que es derivada del denominador:
\[
\int\frac{\dif x}{(x+h)^{2}+s^{2}}=\frac1s\arctan\frac{x+h}{s}+C,\qquad\int\frac{2x+p}{x^{2}+px+q}\dif x=\ln(x^{2}+px+q)+C .
\]

\begin{ejemplo}\label{ej:fracciones}
(a) $\frac{1}{x^{2}-1}=\frac{1}{(x-1)(x+1)}$. Por la \cref{prop:raices-simples}, $A_1=\frac{1}{1-(-1)}=\frac12$ y $A_2=\frac{1}{-1-1}=-\frac12$. Luego $\int\frac{\dif x}{x^{2}-1}=\frac12\ln\Bigl|\frac{x-1}{x+1}\Bigr|+C$.
(b) $\int\frac{x+3}{x^{2}+2x+5}\dif x$. El denominador es $(x+1)^{2}+4$, irreducible. Se escribe $x+3=\frac12(2x+2)+2$:
\[
\int\frac{x+3}{x^{2}+2x+5}\dif x=\frac12\ln(x^{2}+2x+5)+2\cdot\frac12\arctan\frac{x+1}{2}+C .
\]
\end{ejemplo}

\section{Integrales trigonométricas}

Para $\int\sin^{m}x\cos^{n}x\dif x$ con $m,n\in\N_0$ hay dos casos~\cite[§3.2]{openstax2}.
\begin{enumerate}
  \item Si alguno de los exponentes es impar, por ejemplo $m=2k+1$, se aparta un factor $\sin x$, se escribe $\sin^{2k}x=(1-\cos^{2}x)^{k}$ y se sustituye $u=\cos x$. El caso $n$ impar es simétrico, con $u=\sin x$.
  \item Si ambos son pares, se reduce el grado con $\sin^{2}x=\frac{1-\cos2x}{2}$, $\cos^{2}x=\frac{1+\cos2x}{2}$ y $\sin x\cos x=\frac{\sin2x}{2}$.
\end{enumerate}

\begin{ejemplo}\label{ej:trig}
(a) $\int\sin^{3}x\cos^{2}x\dif x=\int(1-\cos^{2}x)\cos^{2}x\,\sin x\dif x$. Con $u=\cos x$, $\dif u=-\sin x\dif x$: $-\int(u^{2}-u^{4})\dif u=\frac{\cos^{5}x}{5}-\frac{\cos^{3}x}{3}+C$.
(b) $\int\sin^{2}x\dif x=\int\frac{1-\cos2x}{2}\dif x=\frac x2-\frac{\sin2x}{4}+C$.
\end{ejemplo}

Las raíces cuadráticas se eliminan con \emph{sustituciones trigonométricas}~\cite[§3.3]{openstax2}: $\sqrt{a^{2}-x^{2}}$ con $x=a\sin t$, como en el \cref{ej:circulo}; $\sqrt{a^{2}+x^{2}}$ con $x=a\tan t$; y $\sqrt{x^{2}-a^{2}}$ con $x=a/\cos t$. Por último, la sustitución $t=\tan\frac x2$ transforma cualquier función racional de $\sin x$ y $\cos x$ en una función racional de $t$, pues $\sin x=\frac{2t}{1+t^{2}}$, $\cos x=\frac{1-t^{2}}{1+t^{2}}$ y $\dif x=\frac{2\dif t}{1+t^{2}}$.

\section{Estrategia}

Ante una integral conviene probar, por este orden: (1) simplificar algebraicamente y buscar una primitiva inmediata (\cref{tab:inmediatas}); (2) buscar una función y su derivada en el integrando, para un cambio de variable; (3) si es un producto de un polinomio o logaritmo por una exponencial o función trigonométrica, integrar por partes; (4) si es racional, dividir y descomponer en fracciones simples; (5) si es trigonométrica o contiene raíces de polinomios de grado~2, usar identidades o sustituciones trigonométricas. Toda primitiva se comprueba derivando.

\begin{observacion}[Errores frecuentes]\label{obs:3-errores}
(a) En un cambio de variable, olvidar transformar $\dif x$: con $u=x^{2}$, $\dif x$ \emph{no} es $\dif u$, sino $\frac{\dif u}{2x}$. (b) Cambiar la variable de una integral definida sin cambiar los límites, o cambiarlos y después deshacer el cambio. (c) Usar un cambio $x=\varphi(t)$ con $\varphi$ discontinua en el intervalo (\cref{eje:3-sustitucion-falsa}). (d) Descomponer en fracciones simples una función racional impropia sin dividir antes. (e) Escribir $\int\frac{\dif x}{x^{2}+1}=\ln(x^{2}+1)$: el numerador no es la derivada del denominador.
\end{observacion}

\begin{observacion}[Conexión con el curso]
Estas técnicas aparecen en todos los cálculos de esperanzas del curso. En el módulo~5, los momentos de las distribuciones uniforme y exponencial se obtienen por partes, y la varianza de la normal también (\cref{tema:4}). En el módulo~18, la discretización de un modelo de espacio de estados contiene $\int_0^{\Delta}e^{a(\Delta-u)}\dif u$, un cambio de variable lineal. En el módulo~21, la esperanza del logaritmo del máximo de variables uniformes se reduce a $\int_0^1u^{n-1}\ln u\dif u=-\frac{1}{n^{2}}$, por partes (\cref{eje:3-curso}). La integral del \cref{eje:3-beta} es la constante de normalización de la distribución beta.
\end{observacion}

\begin{lecturas}
\item Spivak~\cite{spivak}, cap.~19 «Integration in Elementary Terms».
\item OpenStax~\cite{openstax1}, §5.5 «Substitution» y §5.6 «Integrals Involving Exponential and Logarithmic Functions»; \cite{openstax2}, §3.1 «Integration by Parts», §3.2 «Trigonometric Integrals», §3.3 «Trigonometric Substitution», §3.4 «Partial Fractions» y §3.5 «Other Strategies for Integration».
\item Strang~\cite{strang}, §5.4 «Indefinite Integrals and Substitutions» y §7.1--7.4.
\item MIT 18.01SC~\cite{mit1801}, unidad~4, partes~A y~B, sesiones~68--76.
\end{lecturas}

\section*{Ejercicios del tema 3}
\addcontentsline{toc}{section}{Ejercicios del tema 3}

\begin{ejercicio}{1}{eje:3-elegir}
Sin calcularlas, indique qué técnica conviene para cada integral y con qué elección ($u$, $v'$, cambio, descomposición): (a) $\int x\cos(x^{2})\dif x$; (b) $\int x\cos x\dif x$; (c) $\int\frac{\dif x}{x^{2}-5x+6}$; (d) $\int\sin^{2}x\cos^{3}x\dif x$; (e) $\int\sqrt{9-x^{2}}\dif x$; (f) $\int\frac{e^{x}}{1+e^{2x}}\dif x$; (g) $\int x^{3}\ln x\dif x$; (h) $\int\frac{2x+1}{x^{2}+x+7}\dif x$.
\end{ejercicio}

\begin{ejercicio}{2}{eje:3-sustitucion}
Calcule por cambio de variable: (a) $\int\frac{\ln x}{x}\dif x$; (b) $\int x\sqrt{x^{2}-1}\dif x$; (c) $\int x\,e^{-x^{2}}\dif x$; (d) $\int\tan(4x)\dif x$; (e) $\int\frac{x^{2}}{1+x^{6}}\dif x$ (indicación: $u=x^{3}$).
\fuente{adaptado de MIT 18.01SC~\cite{mit1801}, \emph{18.01 Exercises}, sección 5B «Integration by direct substitution».}
\end{ejercicio}

\begin{ejercicio}{2}{eje:3-limites}
Calcule cambiando a la vez la variable y los límites: (a) $\int_1^{e}\frac{\sqrt{\ln x}}{x}\dif x$; (b) $\int_0^{\pi}\frac{\sin x}{(2+\cos x)^{3}}\dif x$; (c) $\int_0^{1}\frac{\dif x}{\sqrt{1-x^{2}}}$ con $x=\sin u$ (es una integral impropia en $x=1$; trátela aquí como límite de $\int_0^{b}$ con $b\to1^{-}$).
\fuente{adaptado de MIT 18.01SC~\cite{mit1801}, \emph{18.01 Exercises}, ejercicio 3E-3.}
\end{ejercicio}

\begin{ejercicio}{2}{eje:3-partes}
Calcule por partes: (a) $\int x\sin x\dif x$; (b) $\int x\,e^{4x}\dif x$; (c) $\int x\ln x\dif x$; (d) $\int\arctan x\dif x$; (e) $\int x^{2}e^{x}\dif x$.
\fuente{adaptado de Strang~\cite{strang}, §7.1, ejercicios 1, 2 y~6.}
\end{ejercicio}

\begin{ejercicio}{2}{eje:3-ciclicas}
Calcule integrando dos veces por partes: (a) $\int e^{x}\cos x\dif x$; (b) $\int\cos(\ln x)\dif x$.
\fuente{adaptado de MIT 18.01SC~\cite{mit1801}, \emph{18.01 Exercises}, ejercicios 5F-4 y 5F-5.}
\end{ejercicio}

\begin{ejercicio}{2}{eje:3-fracciones}
Descomponga en fracciones simples e integre: (a) $\frac{1}{(x-3)(x-2)}$; (b) $\frac{x}{(x-3)(x-2)}$; (c) $\frac{1}{x^{3}-x}$; (d) $\frac{3x+1}{(x-1)^{2}}$.
\fuente{adaptado de Strang~\cite{strang}, §7.4, ejercicios 3, 4, 6 y~8.}
\end{ejercicio}

\begin{ejercicio}{2}{eje:3-fracciones-cuad}
Descomponga en fracciones simples e integre: (a) $\frac{1}{(x-1)(x^{2}+1)}$; (b) $\frac{1}{x^{4}-1}$.
\fuente{adaptado de Strang~\cite{strang}, §7.4, ejercicios 10 y~15.}
\end{ejercicio}

\begin{ejercicio}{2}{eje:3-racional-impropia}
Calcule: (a) $\int\frac{x^{2}+1}{x^{2}+2x+2}\dif x$ (divida primero); (b) $\int\frac{\dif x}{x^{2}+4x+13}$ (complete el cuadrado).
\fuente{adaptado de MIT 18.01SC~\cite{mit1801}, \emph{18.01 Exercises}, ejercicios 5E-10 y 5D (completar el cuadrado).}
\end{ejercicio}

\begin{ejercicio}{2}{eje:3-trig}
Calcule: (a) $\int\sin^{3}x\dif x$; (b) $\int_0^{\pi}\sin^{2}x\dif x$; (c) $\int\cos^{4}x\dif x$ y $\int_0^{\pi/2}\cos^{4}x\dif x$.
\fuente{adaptado de Strang~\cite{strang}, §7.2, ejercicios 1 y~11.}
\end{ejercicio}

\begin{ejercicio}{2}{eje:3-trig-sust}
Calcule con una sustitución trigonométrica: (a) $\int_{-a}^{a}\sqrt{a^{2}-x^{2}}\dif x$ con $a>0$; (b) $\int\frac{x^{2}}{\sqrt{1-x^{2}}}\dif x$; (c) $\int\frac{\dif x}{(1+x^{2})^{2}}$.
\fuente{adaptado de Strang~\cite{strang}, §7.3, ejercicios 3, 5 y~7.}
\end{ejercicio}

\begin{ejercicio}{2}{eje:3-curso}
(a) Sea $a\neq0$ y $\Delta>0$. Calcule $\int_0^{\Delta}e^{a(\Delta-u)}\dif u$ (módulo~18). (b) Sea $n\ge2$ entero. Calcule $\int_0^{1}u^{n-1}\ln u\dif u$ (módulo~21), usando que $u^{n}\ln u\to0$ cuando $u\to0^{+}$.
\end{ejercicio}

\begin{ejercicio}{3}{eje:3-reduccion}
(a) Para $a\neq0$ y $n\in\N$, demuestre la fórmula de reducción $\int x^{n}e^{ax}\dif x=\frac{x^{n}e^{ax}}{a}-\frac{n}{a}\int x^{n-1}e^{ax}\dif x$ y úsela para calcular $\int x^{3}e^{x}\dif x$. (b) Demuestre por inducción sobre $n\in\N_0$ que, para $b>0$,
\[
\int_0^{b}x^{n}e^{-x}\dif x=n!\Bigl(1-e^{-b}\sum_{k=0}^{n}\frac{b^{k}}{k!}\Bigr).
\]
\fuente{adaptado de MIT 18.01SC~\cite{mit1801}, \emph{18.01 Exercises}, ejercicio 5F-2.}
\end{ejercicio}

\begin{ejercicio}{3}{eje:3-wallis}
Sea $I_n=\int_0^{\pi/2}\sin^{n}x\dif x$ para $n\in\N_0$. (a) Demuestre que $I_n=\frac{n-1}{n}I_{n-2}$ para $n\ge2$. (b) Deduzca $I_{2m}=\frac{(2m-1)(2m-3)\cdots1}{(2m)(2m-2)\cdots2}\cdot\frac\pi2$ e $I_{2m+1}=\frac{(2m)(2m-2)\cdots2}{(2m+1)(2m-1)\cdots3}$, y calcule $I_4$ e $I_5$. (c) Demuestre que $n\,I_n\,I_{n-1}=\frac\pi2$ para todo $n\ge1$ y que $(I_n)$ es decreciente.
\end{ejercicio}

\begin{ejercicio}{3}{eje:3-beta}
Demuestre por inducción sobre $n$, integrando por partes, que para todos $m,n\in\N_0$
\[
\int_0^{1}x^{m}(1-x)^{n}\dif x=\frac{m!\,n!}{(m+n+1)!}.
\]
Compruebe la fórmula con $m=3$, $n=2$ desarrollando el integrando.
\end{ejercicio}

\begin{ejercicio}{3}{eje:3-sustitucion-falsa}
Se aplica el cambio $x=1/t$ a $\int_{-1}^{1}\frac{\dif x}{1+x^{2}}$. (a) Compruebe que la sustitución formal produce $-\int_{-1}^{1}\frac{\dif t}{1+t^{2}}$, y que entonces la integral sería igual a su opuesta, luego nula. (b) Calcule el valor correcto. (c) Explique qué hipótesis del \cref{teo:sustitucion} falla, señale en qué paso de su demostración se usa y muestre que el cambio sí es válido en $\int_{1}^{2}\frac{\dif x}{1+x^{2}}$.
\end{ejercicio}
